% ======================================================================
\chapter[Уравнение Шредингера. Бесконечно глубокая яма]%
        {Уравнение Шредингера. Бесконечно глубокая потенциальная яма}
\label{t:06}
% ======================================================================

\begin{qbox}
\textbf{Билет 6.} Уравнение Шредингера. Волновые функции и уровни энергии стационарных
состояний одномерной потенциальной ямы бесконечной глубины.
\end{qbox}

\section{Содержание билета}

\textbf{Подтемы:}
\begin{enumerate}
 \item временно́е уравнение Шредингера: обоснование, свойства (линейность, первый порядок по времени, сохранение нормировки);
 \item стационарные состояния; стационарное уравнение Шредингера как задача на собственные значения $\op H$;
 \item граничные условия для волновой функции;
 \item одномерная бесконечно глубокая яма: уровни энергии, нормированные волновые функции, их свойства;
 \item обобщение на трёхмерный ящик, вырождение.
\end{enumerate}

\textbf{Необходимые понятия:} волновая функция и плотность вероятности (\tref{t:04});
операторы $\op{\vect p}$, $\op H$ и задача на собственные значения (\tref{t:05}); решение
линейного ОДУ с постоянными коэффициентами.

\section{Определения и физические принципы}

\begin{principle}{Уравнение Шредингера (1926)}
Эволюция волновой функции во времени определяется уравнением
\begin{equation}\label{eq:6:tdse}
  \ii\hbar\pdiff{\Psi(\vect r,t)}{t}=\op H\Psi(\vect r,t),\qquad
  \op H=-\frac{\hbar^2}{2m}\nabla^2+U(\vect r,t).
\end{equation}
\end{principle}

\emph{Статус.} Уравнение Шредингера — постулат квантовой механики; его нельзя вывести из
классической физики, но можно обосновать (см. ниже) и проверить опытом. Его роль аналогична
роли уравнений Ньютона в механике или уравнений Максвелла в электродинамике.

\begin{definition}[стационарное состояние]
Состояние с определённым значением энергии $E$. Его волновая функция имеет вид
\begin{equation}\label{eq:6:stat}
  \Psi(\vect r,t)=\psi(\vect r)\,\ee^{-\ii Et/\hbar},
\end{equation}
а $\psi$ удовлетворяет \emph{стационарному уравнению Шредингера}
\begin{equation}\label{eq:6:sse}
  \op H\psi=E\psi,\qquad -\frac{\hbar^2}{2m}\nabla^2\psi+U(\vect r)\psi=E\psi .
\end{equation}
\end{definition}

\emph{Физический смысл.} В стационарном состоянии плотность вероятности $|\Psi|^2=|\psi|^2$ и
средние всех величин, не зависящих явно от времени, постоянны. Именно такие состояния —
«стационарные орбиты» Бора — дают дискретные уровни энергии атомов.

\begin{definition}[потенциальная яма бесконечной глубины]
Одномерная модель частицы, запертой между непроницаемыми стенками:
$U(x)=0$ при $0<x<a$, $U(x)=+\infty$ при $x\le0$ и $x\ge a$.
\end{definition}

\emph{Область применимости модели:} электрон в тонкой плёнке, квантовой яме полупроводниковой
гетероструктуры, в длинной сопряжённой молекуле; нуклон в ядре (грубо). Модель годится,
когда энергия частицы много меньше высоты реальных стенок.

\section{Основная теория}

\subsection{Обоснование уравнения Шредингера}

Требуется линейное уравнение (чтобы выполнялся принцип суперпозиции), первого порядка по
времени (чтобы состояние в момент $t$ определялось только $\Psi(\vect r,0)$), решением которого
для свободной частицы является волна де Бройля $\Psi=A\ee^{\ii(\vect p\vect r-Et)/\hbar}$ с
$E=p^2/2m$. Для неё
\[
  \ii\hbar\pdiff{\Psi}{t}=E\Psi,\qquad -\ii\hbar\nabla\Psi=\vect p\Psi,\qquad
  -\hbar^2\nabla^2\Psi=p^2\Psi .
\]
Соотношение $E=p^2/2m$ превращается в уравнение $\ii\hbar\,\partial_t\Psi=-\frac{\hbar^2}{2m}\nabla^2\Psi$.
Обобщение на частицу в поле получается заменой $E=p^2/2m+U$, что даёт~\eqref{eq:6:tdse}.
Итак, уравнение Шредингера — операторная запись закона сохранения энергии
$E=p^2/2m+U$ с заменами $E\to\ii\hbar\,\partial_t$, $\vect p\to-\ii\hbar\nabla$.

\subsection{Свойства уравнения}

\begin{enumerate}
 \item \textbf{Линейность и однородность} — сумма решений есть решение (принцип суперпозиции).
 \item \textbf{Первый порядок по $t$} — задание $\Psi(\vect r,0)$ однозначно определяет $\Psi(\vect r,t)$:
       квантовая механика детерминистична на уровне волновой функции; случайны лишь результаты измерений.
 \item \textbf{Комплексность} — уравнение содержит $\ii$, поэтому $\Psi$ обязательно комплексна
       (действительное решение не может описывать движение).
 \item \textbf{Сохранение нормировки:} при действительном $U$ выполняется уравнение непрерывности
       $\partial_t|\Psi|^2+\operatorname{div}\vect j=0$ (вывод~\ref{der:6:cont}), и
       $\int|\Psi|^2\dd V$ не зависит от $t$.
\end{enumerate}

\subsection{Стационарные состояния и общее решение}

Если $U$ не зависит от времени, переменные разделяются (вывод~\ref{der:6:sep}), и решения
вида~\eqref{eq:6:stat} существуют для собственных значений $E_n$ оператора $\op H$. Общее
решение — суперпозиция стационарных:
\begin{equation}\label{eq:6:gen}
  \Psi(\vect r,t)=\sum_nc_n\psi_n(\vect r)\,\ee^{-\ii E_nt/\hbar},\qquad
  c_n=\int\psi_n^*(\vect r)\Psi(\vect r,0)\,\dd V .
\end{equation}
Коэффициенты $|c_n|^2$ — вероятности значений энергии $E_n$ — не меняются со временем
(энергия сохраняется), а $|\Psi|^2$ содержит интерференционные члены, осциллирующие с
частотами $(E_n-E_m)/\hbar$ — «боровскими частотами» переходов.

\subsection{Граничные условия}

Стационарное уравнение~\eqref{eq:6:sse} — уравнение второго порядка. Если $U(x)$ имеет конечный
скачок в точке $x_0$, то, интегрируя~\eqref{eq:6:sse} по малому интервалу $(x_0-\varepsilon,x_0+\varepsilon)$, получаем
\[
  \psi'(x_0+\varepsilon)-\psi'(x_0-\varepsilon)=\frac{2m}{\hbar^2}\int_{x_0-\varepsilon}^{x_0+\varepsilon}(U-E)\psi\,\dd x\xrightarrow[\varepsilon\to0]{}0 ,
\]
поскольку подынтегральная функция ограничена. Следовательно, на конечном скачке потенциала
\emph{непрерывны $\psi$ и $\psi'$}. Если же $U=+\infty$ в некоторой области, частица туда
проникнуть не может: $\psi=0$ в этой области, а на её границе $\psi=0$; производная $\psi'$
при этом испытывает скачок (интеграл от $U\psi$ не стремится к нулю).

\subsection{Бесконечно глубокая яма}

Внутри ямы ($0<x<a$) $U=0$, и уравнение~\eqref{eq:6:sse} принимает вид $\psi''+k^2\psi=0$, где
$k^2=2mE/\hbar^2$. С граничными условиями $\psi(0)=\psi(a)=0$ получаем (вывод~\ref{der:6:well})
\begin{equation}\label{eq:6:wellres}
  E_n=\frac{\pi^2\hbar^2}{2ma^2}\,n^2,\qquad
  \psi_n(x)=\sqrt{\frac{2}{a}}\sin\frac{\pi nx}{a},\qquad n=1,2,3,\dots
\end{equation}
Уровни, волновые функции и плотности вероятности для $n=1,2,3$ показаны на рис.~\ref{fig:6:well}.

\begin{figure}[htbp]
\centering
\begin{tikzpicture}
\begin{axis}[width=0.48\linewidth, height=6.6cm, axis lines=none,
  xmin=-0.15, xmax=1.15, ymin=-0.3, ymax=10.8, samples=120, clip=false]
  \addplot[very thick] coordinates {(0,10.5) (0,0) (1,0) (1,10.5)};
  \fill[pattern=north east lines, pattern color=gray] (axis cs:-0.08,0) rectangle (axis cs:0,10.5);
  \fill[pattern=north east lines, pattern color=gray] (axis cs:1,0) rectangle (axis cs:1.08,10.5);
  \foreach \n in {1,2,3} {
     \addplot[gray, dashed, domain=0:1] {\n*\n};
     \addplot[accent, thick, domain=0:1] {\n*\n + 0.85*sin(deg(pi*\n*x))};
  }
  \node[font=\small, anchor=west] at (axis cs:1.1,1) {$E_1$};
  \node[font=\small, anchor=west] at (axis cs:1.1,4) {$E_2=4E_1$};
  \node[font=\small, anchor=west] at (axis cs:1.1,9) {$E_3=9E_1$};
  \node[font=\small] at (axis cs:0,-0.35) {$0$};
  \node[font=\small] at (axis cs:1,-0.35) {$a$};
  \node[font=\small] at (axis cs:0.5,10.3) {$\psi_n(x)$};
\end{axis}
\end{tikzpicture}\hfill
\begin{tikzpicture}
\begin{axis}[width=0.48\linewidth, height=6.6cm, axis lines=none,
  xmin=-0.15, xmax=1.15, ymin=-0.3, ymax=10.8, samples=120, clip=false]
  \addplot[very thick] coordinates {(0,10.5) (0,0) (1,0) (1,10.5)};
  \fill[pattern=north east lines, pattern color=gray] (axis cs:-0.08,0) rectangle (axis cs:0,10.5);
  \fill[pattern=north east lines, pattern color=gray] (axis cs:1,0) rectangle (axis cs:1.08,10.5);
  \foreach \n in {1,2,3} {
     \addplot[gray, dashed, domain=0:1] {\n*\n};
     \addplot[keycol, thick, domain=0:1, fill=keycol, fill opacity=0.12] {\n*\n + 1.5*(sin(deg(pi*\n*x)))^2} \closedcycle;
  }
  \node[font=\small] at (axis cs:0,-0.35) {$0$};
  \node[font=\small] at (axis cs:1,-0.35) {$a$};
  \node[font=\small] at (axis cs:0.5,10.3) {$|\psi_n(x)|^2$};
\end{axis}
\end{tikzpicture}
\caption{Бесконечно глубокая яма: уровни энергии $E_n\propto n^2$ (штриховые линии),
волновые функции $\psi_n$ (слева) и плотности вероятности $|\psi_n|^2$ (справа) для $n=1,2,3$.
Функция $\psi_n$ имеет $n-1$ узел внутри ямы.}
\label{fig:6:well}
\end{figure}

\paragraph{Свойства решения.}
\begin{enumerate}
 \item \textbf{Дискретность} спектра — следствие условия обращения $\psi$ в нуль на обеих стенках:
       на ширине ямы должно укладываться целое число полуволн де Бройля, $a=n\lambda_n/2$.
 \item \textbf{Энергия нулевых колебаний.} Наименьшая энергия $E_1=\pi^2\hbar^2/2ma^2>0$:
       состояние с $E=0$ ($\psi\equiv0$) не нормируемо. Согласие с соотношением
       неопределённостей: $\Delta x\sim a$, $p\gtrsim\hbar/a$, $E\gtrsim\hbar^2/ma^2$.
 \item \textbf{Узлы.} $\psi_n$ имеет $n-1$ нулей внутри ямы (осцилляционная теорема: $n$-е по
       порядку состояние одномерной задачи имеет $n-1$ узлов).
 \item \textbf{Чётность.} Относительно центра ямы $x=a/2$ функции с нечётными $n$ чётны, с
       чётными $n$ — нечётны (следствие симметрии $U(x)$).
 \item \textbf{Ортонормированность:} $\int_0^a\psi_n\psi_m\dd x=\delta_{nm}$.
 \item \textbf{Импульс.} $\psi_n=\frac{1}{\ii\sqrt{2a}}(\ee^{\ii k_nx}-\ee^{-\ii k_nx})$ — суперпозиция волн с
       импульсами $\pm\hbar k_n$; энергия определена, а импульс — нет (его модуль $\hbar k_n$ наиболее
       вероятен, но распределение по $p$ имеет конечную ширину).
\end{enumerate}

\paragraph{Численные оценки.} Для электрона в яме шириной $a=1$~нм:
$E_1=\dfrac{\pi^2(\hbar c)^2}{2\mel c^2a^2}=\dfrac{9{,}87\cdot(197{,}3\ \text{эВ}\cdot\text{нм})^2}{2\cdot0{,}511\cdot10^6\ \text{эВ}\cdot1\ \text{нм}^2}\approx0{,}38$~эВ —
дискретность существенна (сравни $k_{\mathrm B}T\approx0{,}026$~эВ при комнатной температуре).
Для нуклона в ядре ($a\sim10$~фм): $E_1\approx\dfrac{9{,}87\cdot(197{,}3)^2}{2\cdot939\cdot100}\approx2$~МэВ — масштаб
ядерных энергий. Для молекулы газа в сосуде ($a=1$~см) расстояния между уровнями ничтожны — спектр
практически непрерывен.

\paragraph{Трёхмерный ящик.} Для ящика $a\times b\times c$ с непроницаемыми стенками переменные
разделяются, $\psi=\psi_{n_1}(x)\psi_{n_2}(y)\psi_{n_3}(z)$:
\begin{equation}\label{eq:6:box}
  E_{n_1n_2n_3}=\frac{\pi^2\hbar^2}{2m}\Bigl(\frac{n_1^2}{a^2}+\frac{n_2^2}{b^2}+\frac{n_3^2}{c^2}\Bigr),\qquad n_i=1,2,\dots
\end{equation}
Для куба ($a=b=c$) уровни \emph{вырождены}: например, состояния $(2,1,1)$, $(1,2,1)$, $(1,1,2)$
имеют одну энергию $6E_1^{(1D)}$; кратность 3. Вырождение — следствие симметрии куба;
при $a\neq b\neq c$ оно снимается.

\section{Математические выводы}

\begin{derivation}[label=der:6:sep]{разделение переменных; стационарное уравнение}
\emph{Допущение:} $U=U(\vect r)$ не зависит от $t$. Ищем решение~\eqref{eq:6:tdse} в виде
$\Psi(\vect r,t)=\psi(\vect r)f(t)$. Подставляя и деля на $\psi f$:
\[
  \ii\hbar\,\psi\,\diff{f}{t}=f\,\op H\psi
  \quad\Longrightarrow\quad
  \frac{\ii\hbar}{f}\diff{f}{t}=\frac{\op H\psi}{\psi}.
\]
Левая часть зависит только от $t$, правая — только от $\vect r$; значит, обе равны константе $E$:
\[
  \ii\hbar\diff{f}{t}=Ef\ \Rightarrow\ f(t)=\ee^{-\ii Et/\hbar};\qquad \op H\psi=E\psi .
\]
Константа $E$ — энергия: в состоянии $\Psi=\psi\ee^{-\ii Et/\hbar}$ собственное значение оператора
$\op H$ равно $E$, $\Delta E=0$. Плотность вероятности $|\Psi|^2=|\psi|^2|\ee^{-\ii Et/\hbar}|^2=|\psi|^2$
от времени не зависит. Из эрмитовости $\op H$ следует, что $E$ действительна; при комплексном $E$
множитель $|f|^2=\ee^{2\operatorname{Im}E\,t/\hbar}$ нарушал бы сохранение нормировки.
\end{derivation}

\begin{derivation}[label=der:6:cont]{уравнение непрерывности}
Запишем уравнение Шредингера и комплексно сопряжённое к нему ($U$ действительно):
\[
  \ii\hbar\,\partial_t\Psi=-\frac{\hbar^2}{2m}\nabla^2\Psi+U\Psi,\qquad
  -\ii\hbar\,\partial_t\Psi^*=-\frac{\hbar^2}{2m}\nabla^2\Psi^*+U\Psi^* .
\]
Умножим первое на $\Psi^*$, второе на $\Psi$ и вычтем; члены с $U$ сокращаются:
\[
  \ii\hbar\bigl(\Psi^*\partial_t\Psi+\Psi\partial_t\Psi^*\bigr)
  =-\frac{\hbar^2}{2m}\bigl(\Psi^*\nabla^2\Psi-\Psi\nabla^2\Psi^*\bigr).
\]
Слева $\ii\hbar\,\partial_t|\Psi|^2$; справа используем тождество
$\Psi^*\nabla^2\Psi-\Psi\nabla^2\Psi^*=\operatorname{div}(\Psi^*\nabla\Psi-\Psi\nabla\Psi^*)$. Деля на $\ii\hbar$:
\[
  \pdiff{|\Psi|^2}{t}+\operatorname{div}\Bigl[\frac{\hbar}{2m\ii}\bigl(\Psi^*\nabla\Psi-\Psi\nabla\Psi^*\bigr)\Bigr]=0,
  \qquad \vect j=\frac{\hbar}{2m\ii}\bigl(\Psi^*\nabla\Psi-\Psi\nabla\Psi^*\bigr).
\]
Интегрируя по всему пространству и применяя теорему Гаусса (поток через бесконечно удалённую
поверхность равен нулю для нормируемой $\Psi$), получаем $\frac{\dd}{\dd t}\int|\Psi|^2\dd V=0$.
\end{derivation}

\begin{derivation}[label=der:6:well]{уровни и волновые функции бесконечно глубокой ямы}
\emph{Шаг 1 (уравнение).} При $0<x<a$: $-\dfrac{\hbar^2}{2m}\psi''=E\psi$, или
$\psi''+k^2\psi=0$, $k=\sqrt{2mE}/\hbar$. Вне ямы $\psi\equiv0$.

\emph{Шаг 2 (знак энергии).} При $E<0$ решение $\psi=A\ee^{\varkappa x}+B\ee^{-\varkappa x}$ не может обратиться
в нуль в двух точках, кроме тривиального $\psi\equiv0$; при $E=0$ $\psi=A+Bx$ — то же самое.
Следовательно, $E>0$.

\emph{Шаг 3 (общее решение и граничные условия).} $\psi(x)=A\sin kx+B\cos kx$.
\[
  \psi(0)=0\ \Rightarrow\ B=0;\qquad \psi(a)=A\sin ka=0\ \Rightarrow\ ka=\pi n,\ n=1,2,\dots
\]
($A\neq0$, иначе решение тривиально; $n=0$ даёт $\psi\equiv0$; отрицательные $n$ дают те же
функции с обратным знаком, т.\,е. те же состояния).

\emph{Шаг 4 (энергия).}
\[
  k_n=\frac{\pi n}{a}\quad\Longrightarrow\quad E_n=\frac{\hbar^2k_n^2}{2m}=\frac{\pi^2\hbar^2}{2ma^2}n^2 .
\]
\emph{Шаг 5 (нормировка).}
\[
  \int_0^a|A|^2\sin^2\frac{\pi nx}{a}\dd x=|A|^2\int_0^a\frac{1-\cos(2\pi nx/a)}{2}\dd x=|A|^2\frac a2=1
  \ \Rightarrow\ A=\sqrt{\frac2a}.
\]
\emph{Шаг 6 (ортогональность).} При $n\neq m$:
$\frac2a\int_0^a\sin\frac{\pi nx}{a}\sin\frac{\pi mx}{a}\dd x=\frac1a\int_0^a\Bigl[\cos\frac{\pi(n-m)x}{a}-\cos\frac{\pi(n+m)x}{a}\Bigr]\dd x=0$.

\emph{Проверка размерности:} $[\hbar^2/ma^2]=(\text{эрг}\,\text{с})^2/(\text{г}\,\text{см}^2)=\text{эрг}$; $[\psi]=\text{см}^{-1/2}$,
так что $|\psi|^2\dd x$ безразмерна.
\end{derivation}

\begin{derivation}[label=der:6:means]{средние значения в состоянии $\psi_n$}
$\avg x=\frac2a\int_0^ax\sin^2\frac{\pi nx}{a}\dd x=\frac a2$ (симметрия $|\psi_n|^2$ относительно центра).
$\avg{p}=\int\psi_n(-\ii\hbar\psi_n')\dd x=-\ii\hbar\bigl[\psi_n^2/2\bigr]_0^a=0$.
$\avg{p^2}=2mE_n=\pi^2\hbar^2n^2/a^2$, т.\,е. $\Delta p=\pi\hbar n/a$.
$\avg{x^2}=\frac2a\int_0^ax^2\sin^2\frac{\pi nx}{a}\dd x=a^2\Bigl(\frac13-\frac{1}{2\pi^2n^2}\Bigr)$, откуда
\[
  (\Delta x)^2=\avg{x^2}-\frac{a^2}{4}=a^2\Bigl(\frac{1}{12}-\frac{1}{2\pi^2n^2}\Bigr),\qquad
  \Delta x\,\Delta p\big|_{n=1}=\hbar\pi\sqrt{\frac1{12}-\frac1{2\pi^2}}\approx0{,}57\hbar>\frac\hbar2 .
\]
Соотношение неопределённостей выполнено. При $n\to\infty$ $(\Delta x)^2\to a^2/12$ — значение для
равномерного распределения на отрезке, как у классической частицы.
\end{derivation}

\section{Итоговые формулы}

\keyformula{\ii\hbar\pdiff{\Psi}{t}=\op H\Psi,\qquad \op H=-\frac{\hbar^2}{2m}\nabla^2+U}
 {$\Psi(\vect r,t)$ — волновая функция, $\op H$ — гамильтониан.}
 {нерелятивистская частица (без спина).}
 {эволюция состояния во времени по заданному начальному $\Psi(\vect r,0)$.}

\keyformula{\Psi=\psi(\vect r)\ee^{-\ii Et/\hbar},\qquad \op H\psi=E\psi}
 {$E$ — энергия стационарного состояния, $\psi$ — его координатная часть.}
 {$U$ не зависит от времени.}
 {уровни энергии и стационарные волновые функции.}

\keyformula{E_n=\frac{\pi^2\hbar^2}{2ma^2}n^2,\qquad \psi_n=\sqrt{\frac2a}\sin\frac{\pi nx}{a},\quad n=1,2,\dots}
 {$a$ — ширина ямы, $m$ — масса частицы, $n$ — квантовое число.}
 {одномерная яма с бесконечно высокими стенками.}
 {спектр частицы в ограниченной области; энергия нулевых колебаний.}

\keyformula{\psi,\ \psi'\ \text{непрерывны на конечном скачке }U;\qquad \psi=0\ \text{на границе области }U=\infty}
 {$\psi'$ — производная по координате, перпендикулярной границе.}
 {стационарное уравнение Шредингера, кусочно-непрерывный потенциал.}
 {сшивание решений в задачах с кусочно-постоянным потенциалом.}

\section{Физическая интерпретация}

\textbf{Зависимость от параметров.} $E_n\propto n^2/(ma^2)$: чем уже яма и легче частица, тем
больше расстояние между уровнями. Отношение $(E_{n+1}-E_n)/E_n=(2n+1)/n^2\to0$ при $n\to\infty$ —
при больших квантовых числах спектр становится квазинепрерывным (принцип соответствия), а
$|\psi_n|^2$, усреднённая по быстрым осцилляциям, стремится к классическому равномерному
распределению $1/a$.

\textbf{Эксперимент.} Квантование в ямах наблюдается в оптических спектрах полупроводниковых
квантовых ям (GaAs/AlGaAs) и квантовых точек: цвет излучения квантовых точек одного материала
зависит от их размера ($E\propto1/a^2$). Модель «частица в ящике» объясняет также окраску
красителей с длинными сопряжёнными цепями.

\textbf{Сравнение с конечной ямой.} В яме конечной глубины $U_0$ волновая функция проникает под
стенки на глубину $\sim\hbar/\sqrt{2m(U_0-E)}$, эффективная ширина больше $a$, и уровни лежат ниже
значений~\eqref{eq:6:wellres}; число уровней конечно, но хотя бы один уровень в одномерной
симметричной яме существует всегда.

\section{Что необходимо знать на экзамене}

\begin{examcheck}
\checkhead{Определения}
\begin{itemize}
 \item уравнение Шредингера и его свойства; стационарное состояние; граничные условия.
\end{itemize}
\checkhead{Формулы}
\begin{itemize}
 \item $\ii\hbar\,\partial_t\Psi=\op H\Psi$; $\Psi=\psi\ee^{-\ii Et/\hbar}$, $\op H\psi=E\psi$; $\Psi=\sum c_n\psi_n\ee^{-\ii E_nt/\hbar}$;
 \item $E_n=\pi^2\hbar^2n^2/2ma^2$, $\psi_n=\sqrt{2/a}\sin(\pi nx/a)$; трёхмерный ящик и вырождение.
\end{itemize}
\checkhead{Выводы}
\begin{itemize}
 \item «обоснование» уравнения Шредингера из волны де Бройля;
 \item разделение переменных; уравнение непрерывности и сохранение нормировки;
 \item полное решение задачи о бесконечной яме (шаги 1--6), включая нормировку.
\end{itemize}
\checkhead{Интерпретации}
\begin{itemize}
 \item происхождение дискретности (целое число полуволн); $E_1>0$ и соотношение неопределённостей;
 \item узлы и чётность $\psi_n$; переход к классике при $n\to\infty$; численные оценки (электрон в 1~нм — доли эВ).
\end{itemize}
\end{examcheck}
